% Owner: Anshika (Member 2)
%=====================================================================
\section{Conclusion and Future Work}
%=====================================================================
Phase~I showed that separate instruction and data paths remove the structural hazard of a
shared port: CPI\,=\,1.00 against $1+f_{\text{LS}}$. Phase~II moved the comparison to a
realistic hierarchy, with one reusable direct-mapped cache, one shared memory of latency $L$,
and an arbiter. The same RTL and the same memory were measured as a modified Harvard system, as
a single bus with a unified cache of equal capacity, and as a single bus with no cache. An
independent golden model predicted every miss count, and 528 system simulations matched it in
all 9888 checks. Phase~I timing was preserved exactly.

The results lead to three conclusions.
\begin{itemize}
  \item Caches are essential as soon as memory is slower than the core. At $L=8$ the suite CPI
        is 9.88 without a cache, 2.72 with a unified cache and 1.69 with split caches.
  \item The split organisation wins mainly because it keeps two accesses per cycle and keeps
        code and data out of each other's way. With our memory map, most unified misses are
        code/data conflict misses: moving the data by half a cache cuts them from 438 to 165.
        The unified cache wins only when the code footprint is larger than the I-cache, or when
        both caches are tiny.
  \item A one-entry write buffer removes 15\% of the cycles of the split design. It is limited by
        back-to-back stores and by drains that delay instruction refills. Without a cache it
        cannot help at all.
\end{itemize}
The modified Harvard design synthesises without combinational loops, because refill data is
never forwarded and the I-cache request comes from a register.

\textbf{Future work.} The next steps are:
\begin{itemize}
  \item set-associative caches and a victim cache \cite{jouppi1990}, to remove the conflict
        misses measured here;
  \item a write-back D-cache and a deeper write buffer that can be read by loads
        \cite{skadron1997};
  \item burst refills with critical-word-first delivery, which would need a registered bypass
        so as to respect the no-forwarding rule;
  \item a pipelined core, where split caches matter even more because fetch and memory access
        overlap every cycle;
  \item place-and-route on an FPGA board, to measure Fmax and post-route utilisation of the
        cached designs.
\end{itemize}

\section*{Acknowledgment}
The authors thank the faculty and peers of the Modern Computer Architecture laboratory for their
discussions. All performance numbers in this report come from reproducible simulation and
synthesis scripts in the project repository. % TODO (all members): state here how AI tools were
% used, as the course AI policy requires, before submission.
